\documentclass[11pt]{article}
\usepackage[a4paper,margin=27mm]{geometry}
\usepackage{amsmath,amssymb,amsthm}
\usepackage{array}
\usepackage[T1]{fontenc}
\usepackage{lmodern}
\usepackage{microtype}
\usepackage{hyperref}
\hypersetup{colorlinks=true,linkcolor=blue,citecolor=blue,urlcolor=blue,
pdftitle={OEIS A391810: formula (5), COMMENTS (C.1)-(C.10), and prime-index vanishing},
pdfauthor={Haobo Ma; Wenlin Zhang}}
\newtheorem{theorem}{Theorem}
\newtheorem{lemma}{Lemma}
\newtheorem{corollary}{Corollary}
\newcommand{\Z}{\mathbb Z}
\newcommand{\Q}{\mathbb Q}
\newcommand{\round}{\operatorname{round}}
\title{OEIS A391810: formula (5),\\
COMMENTS (C.1)--(C.10), and prime-index vanishing}
\author{Haobo Ma\\
\small ChronoAI Pte Ltd, Singapore
\and Wenlin Zhang\\
\small National University of Singapore\\
\small ChronoAI Pte Ltd, Singapore}
\date{2 November 2026}
\begin{document}
\maketitle
\begin{center}\small
Licensed under Creative Commons Attribution 4.0 International (CC BY 4.0).
\end{center}
\begin{abstract}
We prove formula (5), COMMENTS (C.1)--(C.10), and the zero clause of
COMMENT (C.11) of Paul D. Hanna's OEIS entry A391810,
revision 39~\cite{hanna}.
For each positive integer $n$, let $T(n,k)$ be the coefficient of $x^{-k}$
in the formal bilateral Laurent series
$\sum_{j\in\Z}x^j(1-x^j)^{n+j}$ at zero.
We prove that the negative alternating sum of the complete row is
$2\round(2^n/3)=2J_n$, where $J_n$ is the offset-zero Jacobsthal sequence.
The proof identifies every nonpositive Laurent coefficient with a coefficient
of a finite polynomial, proves the stated row bound, and evaluates that
polynomial at $-1$. No substitution into the bilateral series is required.
The same polynomial proves the conjectured row sums, first and second
moments, and alternating first-moment generating function in formulas
(2), (3), (4), and (6) of that entry. A divisor-sum expression gives
the individual terms and shows that $T(n,q)=0$ for every prime $q>n$.
The further assertion in (C.11) that the next prime is the first zero
remains open here.
\end{abstract}

\section{The identity and its conventions}

Formula (5) and COMMENTS (C.1)--(C.11) of Paul D. Hanna's OEIS entry A391810
pose alternating-sum and individual-term conjectures~\cite{hanna}.
The negative
alternating sum of the full signed Laurent row equals twice a Jacobsthal
number. We establish this identity and the ten individual-term formulas,
together with prime-index vanishing, using the original row
interpretation and inclusive endpoint. The Jacobsthal closed form and
nearest-integer description are already recorded in A001045~\cite{jacobsthal};
the issue here is their connection to the bilateral Laurent coefficients.
The pinned formal source~\cite{formal} supplies the same connection and
the whole formula. This note gives a self-contained written proof.

For $n\geq1$, work in the formal Laurent field $\Q((x))$ and put
\begin{equation}\label{eq:series}
 S_n(x)=\sum_{j\in\Z}x^j(1-x^j)^{n+j},\qquad
 T(n,k)=[x^{-k}]S_n(x),\quad k\geq0.
\end{equation}
Integer powers with negative exponents mean inverses in this field.
The $j=0$ term is zero, since $n>0$.
We justify the coefficientwise sum below; it is not an analytic convergence
assertion. Write
\[
 B_n=\left\lfloor\frac{n+1}{2}\right\rfloor^2,
 \qquad J_0=0,\quad J_1=1,\quad J_{n+2}=J_{n+1}+2J_n.
\]
For rational arguments, $\round$ denotes nearest-integer rounding, with ties
toward positive infinity. The arguments used here are never ties.

\begin{theorem}\label{thm:main}
For every integer $n\geq1$, all $T(n,k)$ are integers, $T(n,k)=0$ for
$k>B_n$, and
\begin{equation}\label{eq:main}
 -\sum_{k=0}^{B_n}(-1)^k T(n,k)
 =2\round\!\left(\frac{2^n}{3}\right)=2J_n.
\end{equation}
\end{theorem}

\section{From the bilateral series to a finite row}

\begin{lemma}\label{lem:row}
For every $n\geq1$ and $k\geq0$,
\begin{equation}\label{eq:bridge}
 T(n,k)=[y^k]P_n(y),\qquad
 P_n(y)=-1+\sum_{h=1}^{n}y^h(1-y^h)^{n-h}\in\Z[y].
\end{equation}
Moreover, $\deg P_n\leq B_n$.
\end{lemma}

\begin{proof}
First consider the formal sum in~\eqref{eq:series}. For $j>0$, its $j$th
summand has least exponent $j$. For $j=-h<0$, the identity
$1-x^{-h}=-x^{-h}(1-x^h)$ gives
\begin{equation}\label{eq:normal}
 x^{-h}(1-x^{-h})^{n-h}
 =(-1)^{n-h}x^{h(h-n-1)}(1-x^h)^{n-h}.
\end{equation}
The last factor has constant coefficient one and only nonnegative
exponents, also when $n-h<0$. Thus the least exponent is $h(h-n-1)$.
These least exponents tend to positive infinity as $h\to\infty$;
the positive-index least exponents do likewise. They are uniformly
bounded below, and only finitely many summands can contribute to any fixed
coefficient. This proves that the bilateral coefficientwise sum is a
well-defined element of $\Q((x))$.

No $j>0$ term contributes to a nonpositive power of $x$. For $1\leq h\leq n$,
the $j=-h$ term is the finite polynomial
$y^h(1-y^h)^{n-h}$ in $y=x^{-1}$.
At the boundary $h=n+1$, it is
\[
 x^{-(n+1)}(1-x^{-(n+1)})^{-1}
 =-\frac{1}{1-x^{n+1}}
 =-\sum_{r\geq0}x^{(n+1)r}.
\]
Its nonpositive part is exactly the constant $-1$.
For $h\geq n+2$, the least exponent $h(h-n-1)$ is positive.
This proves~\eqref{eq:bridge}, including the constant coefficient.

It remains to check the sharper row bound. The $h$th summand of $P_n+1$
has degree $h(n+1-h)$. If $n=2m+1$, this is at most $(m+1)^2=B_n$,
as follows by completing the square.
If $n=2m$, split off its highest monomial
\[
 (-1)^{2m-h}y^{h(2m+1-h)},\qquad 1\leq h\leq2m.
\]
Pairing $h$ with $2m+1-h$ cancels these monomials: their exponents
coincide and their signs are opposite. There is no fixed index in this
pairing. Every remaining monomial in the $h$th summand has exponent
$h(r+1)$ with $0\leq r\leq2m-h-1$, hence at most
$h(2m-h)\leq m^2=B_n$. When that range is empty, there is no remaining
monomial. The constant term also respects the bound. This proves the lemma.
\end{proof}

\section{Polynomial evaluation and Jacobsthal numbers}

\begin{proof}[Proof of Theorem~\ref{thm:main}]
By Lemma~\ref{lem:row}, the full sum in~\eqref{eq:main} is $-P_n(-1)$.
We evaluate the finite polynomial, not $S_n$.
An odd index $h$ contributes $-2^{n-h}$ to $P_n(-1)$.
An even index $h<n$ contributes zero; if $h=n$ is even, its contribution
is one, since its exponent $n-h$ is zero.
For $n=2m+1$ with $m\geq0$, the odd indices give
\[
 P_{2m+1}(-1)=-1-\sum_{r=0}^{m}4^r
 =-\frac{4^{m+1}+2}{3}
 =-\frac{2(2^{2m+1}+1)}{3}.
\]
For $n=2m$ with $m\geq1$, the even boundary term cancels the constant
$-1$, so
\[
 P_{2m}(-1)=-2\sum_{r=0}^{m-1}4^r
 =-\frac{2(4^m-1)}{3}.
\]
\par\nopagebreak
Together these yield
\begin{equation}\label{eq:value}
 -P_n(-1)=\frac{2\bigl(2^n-(-1)^n\bigr)}{3}.
\end{equation}

To identify the right-hand side, put $C_n=(2^n-(-1)^n)/3$.
Then $C_0=0$, $C_1=1$, and a direct calculation gives
\[
 C_{n+1}+2C_n
 =\frac{2^{n+2}-\bigl((-1)^{n+1}+2(-1)^n\bigr)}{3}
 =\frac{2^{n+2}-(-1)^{n+2}}{3}=C_{n+2}.
\]
Induction in the defining recurrence therefore gives $C_n=J_n$ for all
$n\geq0$. In particular,
\[
 \frac{2^n}{3}-J_n=\frac{(-1)^n}{3}.
\]
The absolute error is $1/3<1/2$, so the unique nearest integer to $2^n/3$
is $J_n$. Combining this with~\eqref{eq:value} proves both equalities.
Integrality and the vanishing assertion follow from Lemma~\ref{lem:row}.
\end{proof}

The first six alternating sums are $2,2,6,10,22,42$.
The positive-index restriction matters: the finite expression gives
$P_0(y)=-1$, so its negative value at $-1$ is $1$, whereas $2J_0=0$.
No zero-index row extension is asserted.

\section{Individual terms and prime-index zeros}

We use $\binom{a}{b}=0$ when $a\geq0$ and $b$ is outside $0\leq b\leq a$.

\begin{corollary}[Divisor formula]\label{cor:divisors}
For $n\geq1$, $T(n,0)=-1$, and for every $k\geq1$,
\begin{equation}\label{eq:divisors}
 T(n,k)=\sum_{\substack{h\mid k\\1\leq h\leq n}}
 (-1)^{k/h-1}\binom{n-h}{k/h-1}.
\end{equation}
\end{corollary}

\begin{proof}
In the $h$th summand of $P_n+1$, the binomial expansion is
\[
 y^h(1-y^h)^{n-h}
 =\sum_{r=0}^{n-h}(-1)^r\binom{n-h}{r}y^{h(r+1)}.
\]
Its $y^k$ coefficient is zero unless $h\mid k$; otherwise put
$r=k/h-1$. All these summands have positive degree, so the constant
coefficient of $P_n$ is $-1$.
\end{proof}

\begin{corollary}[COMMENTS (C.1)--(C.5)]\label{cor:small}
The first five individual-term conjectures hold in their stated ranges:
\begin{align*}
 T(n,0)&=-1 &&(n\geq1),&
 T(n,1)&=1 &&(n\geq1),\\
 T(n,2)&=-(n-2) &&(n\geq2),&
 T(n,3)&=\frac{n(n-3)}2+2 &&(n\geq3),\\
 T(n,4)&=-\frac{(n-3)(n^2-3n+8)}6 &&(n\geq4).
\end{align*}
\end{corollary}

\begin{proof}
The constant term is given above. For $k=1,2,3,4$, respectively,
formula~\eqref{eq:divisors} gives
\[
 1,\qquad -\binom{n-1}{1}+1,\qquad
 \binom{n-1}{2}+1,\qquad
 -\binom{n-1}{3}-\binom{n-2}{1}+1.
\]
The divisors involved are $\{1\}$, $\{1,2\}$, $\{1,3\}$, and
$\{1,2,4\}$. Expanding the displayed binomials gives the stated formulas.
\end{proof}

\begin{corollary}[COMMENTS (C.6)--(C.9)]\label{cor:edges}
For $m\geq1$,
\[
 T(2m-1,m^2)=(-1)^{m-1},\qquad
 T(2m,m^2)=(-1)^{m-1}m.
\]
For $m\geq2$,
\[
 T(2m-1,m^2-1)=2(-1)^m,\qquad
 T(2m,m^2-1)=2m(-1)^m.
\]
\end{corollary}

\begin{proof}
For the odd row, the $h$th summand has degree
$h(2m-h)=m^2-(h-m)^2$. Only $h=m$ reaches $m^2$;
its highest coefficient is $(-1)^{m-1}$. To reach $m^2-1$, only
$h=m-1,m,m+1$ are possible. The middle index contributes nothing,
because $m$ does not divide $m^2-1$ when $m\geq2$.
The other two contribute their highest coefficients, both $(-1)^m$.

For the even row, first remove the highest monomial of each summand.
Their total is zero by the pairing used in Lemma~\ref{lem:row}.
The remaining $h$th polynomial has degree at most
$h(2m-h)=m^2-(h-m)^2$; when nonempty, its highest coefficient is
\[
 (-1)^{2m-h-1}\binom{2m-h}{2m-h-1}
 =(-1)^{2m-h-1}(2m-h).
\]
Thus $h=m$ alone contributes to $m^2$, with coefficient
$(-1)^{m-1}m$. For $m^2-1$ and $m\geq2$, only $h=m-1,m+1$
contribute, since the same divisibility argument excludes $h=m$.
Their contributions are $(-1)^m(m+1)$ and $(-1)^m(m-1)$.
\end{proof}

\begin{corollary}[COMMENT (C.10) and the zero clause of (C.11)]\label{cor:primes}
For every odd prime $p$, $T(p,p)=2$. For every $n\geq1$ and prime
$q>n$, $T(n,q)=0$. In particular, the next prime after $n$ has zero
coefficient, as asserted for $n\geq5$ in (C.11).
\end{corollary}

\begin{proof}
For $T(p,p)$ the divisors in~\eqref{eq:divisors} are $1,p$.
Since $p$ is odd their contributions are
\[
 (-1)^{p-1}\binom{p-1}{p-1}=1
 \qquad\text{and}\qquad \binom00=1.
\]
If $q>n$, only the divisor $1$ is allowed, and
$\binom{n-1}{q-1}=0$.
\end{proof}

The stronger assertion in (C.11) is that $T(n,k)\ne0$ for every
$0\leq k<q$, where $q$ is the next prime after $n\geq5$.
That first-zero assertion is not established by the divisor formula;
it remains an open question in this note. Exact checks of rows $5\leq n\leq79$
agree with it.

\section{The sibling conjectures}

All sums below retain the complete inclusive endpoint $B_n$.
Let $D=y\,d/dy$ denote the Euler differential operator.

\begin{corollary}\label{cor:moments}
Formulas (2) and (3) of A391810 hold for every $n\geq1$:
\[
 \sum_{k=0}^{B_n}T(n,k)=0,\qquad
 \sum_{k=0}^{B_n}kT(n,k)=1.
\]
Formula (4) holds for every $n\geq2$:
\[
 \sum_{k=0}^{B_n}k^2T(n,k)=5-2n.
\]
\end{corollary}

\begin{proof}
At $y=1$, only the $h=n$ summand of $P_n+1$ survives, giving
$P_n(1)=0$. In $P_n'(1)$, the terms $h=n,n-1$ give
$n-(n-1)=1$; for $n=1$, use $P_1(y)=-1+y$ directly.
For $D^2P_n(1)$, only $n-h=0,1,2$ survive. Applying $D^2y^a=a^2y^a$
to $y^h$, $y^h-y^{2h}$ and $y^h-2y^{2h}+y^{3h}$ gives
\[
 (D^2P_n)(1)=n^2-3(n-1)^2+2(n-2)^2=5-2n.
\]
When $n=2$, the absent $h=0$ term has the displayed weight zero.
Lemma~\ref{lem:row} identifies these three evaluations with the stated sums.
\end{proof}

\begin{corollary}\label{cor:alternating-moment}
Formula (6) of A391810 holds. Namely, for
\[
 A_n=-\sum_{k=0}^{B_n}(-1)^k kT(n,k)=P_n'(-1),\qquad n\geq1,
\]
the formal ordinary generating function is
\begin{equation}\label{eq:moment-gf}
 \sum_{n\geq1}A_nz^n
 =\frac{z(1-z+8z^2)}{(1+z)^2(1-2z)^2}.
\end{equation}
\end{corollary}

\begin{proof}
Write $n=h+r$ and $f_{h,r}(y)=y^h(1-y^h)^r$.
For odd $h$, direct differentiation gives
$f_{h,r}'(-1)=h(2^r+r2^{r-1})$, with the second term zero when $r=0$.
For even $h$, the derivatives for $r=0,1$ are $-h,h$, and all others vanish.
Consequently,
\[
 \sum_{r\geq0}z^{h+r}f_{h,r}'(-1)=
 \begin{cases}
 hz^h(1-z)/(1-2z)^2,&h\text{ odd},\\
 hz^h(z-1),&h\text{ even}.
 \end{cases}
\]
Each coefficient involves finitely many $h$, so rearrangement in $\Q[[z]]$
is valid. Differentiating the geometric series gives
\[
 \sum_{h\text{ odd}}hz^h=\frac{z(1+z^2)}{(1-z^2)^2},\qquad
 \sum_{h\text{ even}}hz^h=\frac{2z^2}{(1-z^2)^2}.
\]
These identities yield
\[
 \sum_{n\geq1}A_nz^n
 =\frac{1-z}{(1-z^2)^2}
   \left(\frac{z(1+z^2)}{(1-2z)^2}-2z^2\right)
 =\frac{z(1-z+8z^2)}{(1+z)^2(1-2z)^2}.
\]
\end{proof}

The first eight $A_n$ are $1,1,13,25,81,181,453,1025$, as in the source.
Formula (1) defines the row. Formula (7), which is not labeled a conjecture
there, is the defining squared-norm relation for the companion entry
A391811~\cite{norm}. Its half-sum of squares is an integer:
$T(n,k)^2\equiv T(n,k)\pmod2$ and the row sum is zero.
It is positive, since $T(n,0)=-1$ and $T(n,1)=1$.
The row-sum COMMENT in A391811 is also Corollary~\ref{cor:moments}.

\paragraph{Coverage of the posing entry.}
The source locators in the following table refer to revision 39 of A391810.
\begin{center}\small
\begin{tabular}{>{\raggedright\arraybackslash}p{0.36\linewidth}>{\raggedright\arraybackslash}p{0.55\linewidth}}
\hline
Source item & Mathematical coverage\\ \hline
Formula (1); row endpoint & Definition; Lemma~\ref{lem:row}\\
Formulas (2), (3), (4) & Corollary~\ref{cor:moments}\\
Formula (5) & Theorem~\ref{thm:main}\\
Formula (6) & Corollary~\ref{cor:alternating-moment}\\
COMMENTS (C.1)--(C.5) & Corollary~\ref{cor:small}\\
COMMENTS (C.6)--(C.9) & Corollary~\ref{cor:edges}\\
COMMENT (C.10) & Corollary~\ref{cor:primes}\\
(C.11), next-prime coefficient zero & Corollary~\ref{cor:primes}, for every prime $q>n$\\
(C.11), first-zero location & Open here; finite checks for $5\leq n\leq79$\\
``Each row sum appears to equal zero'' & Same statement as formula (2)\\
Formula (7); A391811 & Defining relation; positive integrality as above\\
Two related bilateral identities in COMMENTS & Outside the nonpositive-coefficient results here\\ \hline
\end{tabular}
\end{center}

\paragraph{Source and verification scope.}
The result is stated as
\href{https://github.com/the-omega-institute/trureturing/blob/1239bd2a0ed6e9d54aabb5a42af94234e9eb61ae/D5/S3/ArithSums/HannaLaurentAlternatingJacobsthal.lean#L103}{\texttt{full\_formula\_five}} in the pinned
module in~\cite{formal}. Its five clauses retain the integer
row sum, the rounding--Jacobsthal equality, the rational original-Laurent sum,
the coefficient bridge for every $k\geq0$, and vanishing beyond $B_n$.
The companion
\href{https://github.com/the-omega-institute/trureturing/blob/1239bd2a0ed6e9d54aabb5a42af94234e9eb61ae/D5/S3/ArithSums/HannaLaurentRowSums.lean}{\texttt{HannaLaurentRowSums.lean}}
contains \texttt{series\_coefficient}, \texttt{row\_degree}, and
\texttt{result}: the coefficient bridge, support bound, and formula (2).
The declarations \texttt{result} in
\href{https://github.com/the-omega-institute/trureturing/blob/1239bd2a0ed6e9d54aabb5a42af94234e9eb61ae/D5/S3/ArithSums/HannaLaurentFirstMoment.lean}{\texttt{HannaLaurentFirstMoment.lean}}
and
\href{https://github.com/the-omega-institute/trureturing/blob/1239bd2a0ed6e9d54aabb5a42af94234e9eb61ae/D5/S3/ArithSums/HannaLaurentSecondMoment.lean}{\texttt{HannaLaurentSecondMoment.lean}}
state formulas (3) and (4). All four modules have the revision specified
in~\cite{formal}. The divisor formula, COMMENTS (C.1)--(C.10),
prime-index vanishing, and formula (6) have the written proofs above.

The companion script \texttt{verify.py}, with results in
\texttt{verification.json}, compares the divisor formula and an independent
Laurent expansion with the polynomial on full rows $1\leq n\leq79$,
checks the term formulas through row 79, and checks all four moment identities.
The Jacobsthal scalar identity is checked through $n=512$.
These are finite exact-arithmetic checks supplementing the written proofs.

\begin{thebibliography}{9}
\raggedright
\bibitem{hanna}
P. D. Hanna, \emph{Entry A391810}, The On-Line Encyclopedia of Integer Sequences,
revision 39, 3 February 2026; author line 21 December 2025.
Definition, row endpoint, formulas (2)--(6), and COMMENTS (C.1)--(C.11).
\url{https://oeis.org/A391810}.

\bibitem{jacobsthal}
The OEIS Foundation Inc., \emph{A001045: Jacobsthal sequence},
The On-Line Encyclopedia of Integer Sequences,
revision 1350, 30 June 2026.
Name, offset, recurrence, and closed-form/rounding formulas.
\url{https://oeis.org/A001045}.

\bibitem{norm}
P. D. Hanna, \emph{Entry A391811}, The On-Line Encyclopedia of Integer Sequences,
revision 11, 28 December 2025; author line 21 December 2025.
Definition and formula for the half-sum of squared A391810 row entries;
row-sum COMMENT.
\url{https://oeis.org/A391811}.

\bibitem{formal}
Trureturing contributors,
\emph{Formal sources for A391810}, directory \texttt{D5/S3/ArithSums}, revision
\nolinkurl{1239bd2a0ed6e9d54aabb5a42af94234e9eb61ae}.
\href{https://github.com/the-omega-institute/trureturing/blob/1239bd2a0ed6e9d54aabb5a42af94234e9eb61ae/D5/S3/ArithSums/HannaLaurentAlternatingJacobsthal.lean}{\texttt{HannaLaurentAlternatingJacobsthal.lean}}, theorem
\texttt{full\_formula\_five}; companion modules linked above.
\end{thebibliography}
\end{document}
